% ============================================================
%  chapters/07_conclusion.tex
% ============================================================

\section{Conclusion}
\label{sec:conclusion}

The purpose of this paper was to determine whether a link exists between
community-level violence and educational attainment and school enrollment in
Nigeria. Using the 2010 NEDS household survey and the UCDP violence dataset
for 1999--2010, this analysis does not find a clear, consistent causal link
between violence in a child's LGA and household education consumption
decisions. However, two important patterns emerge.

First, the analysis confirms several well-established findings from the
broader NEDS survey. Men achieve 1.5 more years of schooling than women in
urban areas and 1.8 more years in rural areas, and primary school attendance
is 64\% for boys versus 58\% for girls \citep{neds2010}. These gender gaps
are confirmed in our regressions, where the coefficient on \textit{male} is
positive and statistically significant across all specifications. Urban
children are also consistently more likely to attend primary school (74\% vs.
55\%) and secondary school (60\% vs. 36\%), a pattern borne out by the
positive and generally significant coefficient on \textit{urban}.

Second, the reversal in sign of the violence coefficients when moving from
the baseline OLS to the LGA fixed effects specification---particularly for
the \textit{highest\_enrolled} outcome---highlights a fundamental
identification challenge. Violence is not randomly assigned across LGAs. In
Nigeria, LGAs that experienced violence during 1999--2010 differ
systematically from those that did not in ways that correlate with
educational outcomes. The OLS results absorb some of this selection through
covariates but likely not all of it. The fixed effects estimates eliminate
time-invariant LGA-level confounders but cannot address potential
simultaneity or time-varying confounders. Cleaner causal identification---
perhaps through an instrumental variables strategy or a difference-in-
differences design exploiting geographic proximity to conflict sites---would
be an important avenue for future work.

\subsection{Policy Implications}

Although a clear causal link between violence and educational attainment was
not established, the results do point to several actionable policy
directions.

The persistent urban-rural gap in school attendance, driven in large part by
distance to school, calls for expanded geographic access to schooling in
rural communities. Survey data indicate that 85\% of urban children live
within 15 minutes of their school, compared with 62\% of rural children. In
the northern region, average walk times to school range from 23 to 37
minutes, versus 14 to 19 minutes in the South \citep{neds2010}. Distance
was reported by parents as one of the primary reasons children started school
late or did not attend at all. Targeted school construction programs in
underserved areas could meaningfully reduce this barrier.

Monetary costs also remain a binding constraint. Even though the UBE scheme
makes primary and junior secondary school free, families must pay for
uniforms, books, and supplies out of pocket. Better management of school
funds and potentially subsidized supply provision could reduce this burden
for the poorest households.

\subsection{Limitations and Future Research}

The most important limitation of this analysis is the absence of a credible
identification strategy for isolating the causal effect of violence.
Future work should pursue one of the following approaches:

\begin{enumerate}
    \item \textbf{Difference-in-differences} exploiting the onset of the
          Boko Haram insurgency in 2009 as an exogenous shock, comparing
          pre- and post-2009 educational outcomes across high- and low-
          exposure LGAs using multiple DHS rounds (2008, 2013, 2018).
    \item \textbf{Instrumental variables} using geographic or conflict-
          specific variation in violence exposure that is plausibly
          exogenous to educational outcomes.
    \item \textbf{Richer data} with complete and reliable violence records
          at the LGA level and individual-level panel data to track
          educational trajectories over time.
\end{enumerate}

The data limitations encountered here---particularly missing values,
imputed variables, and erroneous data entry---also point to the importance
of investing in higher-quality administrative and survey data collection
infrastructure in Nigeria. Without reliable data, researchers cannot reach
strong conclusions, and policymakers cannot design effective interventions.
